\chapter{Bloc 2 --- Collecte \& Stockage (C8-C12)}
\label{chap:bloc2}

Le Bloc 2 industrialise la collecte multi-source, nettoie et consolide
les données dans une base de travail, puis les expose via une API.
Clos le 27/08/2026 (milestone M1), les 5 compétences ont chacune un
livrable vérifié -- détail exhaustif dans
\texttt{docs/comptes\_rendus/M1.md}.

\section{Séquencement retenu}

L'ordre C8~$\rightarrow$~C10~$\rightarrow$~C9~$\rightarrow$~C11~$\rightarrow$~C12
(\texttt{docs/architecture/sequencement\_bloc2.md}), documenté et
justifié avant tout développement plutôt que décidé au fil de l'eau,
diffère de l'ordre naturel de lecture du cahier des charges (qui aurait
suggéré de modéliser avant d'agréger). L'expérience a validé ce choix :
le grain réel des doublons dans les fichiers clients n'a été découvert
qu'en les traitant concrètement, pas en les modélisant par anticipation.

\section{C8 --- Automatiser l'extraction de données}

Cinq connecteurs couvrent l'ensemble des sources identifiées en C2 :
API REST (TransFlow), scraping du portail transporteur, requêtes SQL
directes (FluxPro), fichiers clients hétérogènes, et l'historique
volumineux -- testés de bout en bout (31 tests).

\section{C9 --- Développer des requêtes SQL d'extraction}

Cinq requêtes documentées avec leurs résultats réels
(\texttt{docs/architecture/requetes\_sql\_extraction.md}). L'historique
est chargé dans la base de staging PostgreSQL existante plutôt que dans
un outil « big data » séparé : au volume réel du jeu de données
(25~000 lignes), un outil dédié aurait été disproportionné --- choix
justifié par la sobriété numérique (RGESN) plutôt que par principe.

\section{C10 --- Développer des règles d'agrégation de données}

Le nettoyage et le dédoublonnage des fichiers clients sont menés au bon
grain (commande, produit) : 13 à 17~\% de doublons réels mesurés sur le
jeu de données, pas un chiffre théorique.

\section{C11 --- Créer une base de données}

La base de travail est modélisée selon la méthode MERISE
(\texttt{docs/architecture/modelisation\_merise.md}), versionnée via
Alembic. L'analyse des formes normales a été menée empiriquement --- en
interrogeant réellement les données --- plutôt que théoriquement, et a
révélé deux violations réelles, corrigées après vérification chiffrée :
\texttt{livraisons.statut} (3NF, remplacé par une vue SQL) et
\texttt{commandes\_clients} (2NF, décomposée en en-tête/lignes).

La clé métier de rapprochement entre livraisons et expéditions FluxPro
est \texttt{tracking\_number}, portée sans contrainte de clé étrangère
formelle (les deux tables appartiennent à des systèmes sources
distincts) :

\begin{lstlisting}[language=Python,caption={Clé de rapprochement \texttt{tracking\_number} (\texttt{src/datacore/storage/staging/models.py})}]
# tracking_number : cle metier de rapprochement avec
# expeditions.tracking_number (FluxPro) -- pas de contrainte FK
# (systemes sources distincts), rapprochement fait a la lecture.
Column("tracking_number", String(20), nullable=False),
\end{lstlisting}

\section{C12 --- Partager le jeu de données}

L'API Omega Data (FastAPI, choisie pour la génération automatique de
la spécification OpenAPI exigée par le livrable) expose le jeu de
données consolidé, avec authentification par clé et autorisation par
rôle -- les référents clients externes sont systématiquement
restreints à leur propre client, règle reprise directement du registre
RGPD établi en C11.

\section{Point ouvert transmis au bloc suivant}

Le rapprochement entre \texttt{clients.nom} (FluxPro) et le champ texte
libre \texttt{historique\_expeditions.client} reste non garanti par une
contrainte formelle à ce stade --- signalé comme point à traiter
explicitement lors de la modélisation de l'entrepôt de données
(chapitre~\ref{chap:bloc3}).
